% paper08-treaty-synthesis.tex — Treaty Synthesis in Hypercover Families
% Standalone paper for the Judgment Geometry project
\documentclass[11pt]{article}
% jugeo-common.tex — shared preamble for all Judgment Geometry papers
% Include via: % jugeo-common.tex — shared preamble for all Judgment Geometry papers
% Include via: % jugeo-common.tex — shared preamble for all Judgment Geometry papers
% Include via: \input{jugeo-common}

% ── Packages ────────────────────────────────────────────────────────────
\usepackage{amsmath,amssymb,amsthm,mathtools}
\usepackage{tikz-cd}
\usepackage{listings}
\usepackage{xcolor}
\usepackage{xspace}
\usepackage{booktabs}
\usepackage{multirow}
\usepackage{enumitem}
\usepackage[expansion=false]{microtype}
\usepackage{hyperref}
\usepackage{cleveref}
% \usepackage{thmtools}  % disabled: not available in this TeX installation
\usepackage{float}
\usepackage{caption}
\usepackage{subcaption}
\usepackage{tabularx}
\usepackage{stmaryrd}       % \llbracket, \rrbracket
\usepackage{bbm}            % \mathbbm{1}

% ── Page geometry ───────────────────────────────────────────────────────
\usepackage[margin=1in]{geometry}

% ── Theorem environments ────────────────────────────────────────────────
\theoremstyle{plain}
\newtheorem{theorem}{Theorem}[section]
\newtheorem{lemma}[theorem]{Lemma}
\newtheorem{proposition}[theorem]{Proposition}
\newtheorem{corollary}[theorem]{Corollary}

\theoremstyle{definition}
\newtheorem{definition}[theorem]{Definition}
\newtheorem{example}[theorem]{Example}
\newtheorem{construction}[theorem]{Construction}

\theoremstyle{remark}
\newtheorem{remark}[theorem]{Remark}
\newtheorem{notation}[theorem]{Notation}

% ── Python listings style ──────────────────────────────────────────────
\definecolor{jugeoBlue}{HTML}{2563EB}
\definecolor{jugeoGreen}{HTML}{059669}
\definecolor{jugeoRed}{HTML}{DC2626}
\definecolor{jugeoPurple}{HTML}{7C3AED}
\definecolor{jugeoGray}{HTML}{6B7280}
\definecolor{jugeoBg}{HTML}{F8FAFC}

\lstdefinestyle{jugeo-python}{
  language=Python,
  basicstyle=\ttfamily\small,
  keywordstyle=\color{jugeoBlue}\bfseries,
  stringstyle=\color{jugeoGreen},
  commentstyle=\color{jugeoGray}\itshape,
  numberstyle=\tiny\color{jugeoGray},
  backgroundcolor=\color{jugeoBg},
  numbers=left,
  numbersep=6pt,
  frame=single,
  framerule=0.4pt,
  rulecolor=\color{jugeoGray!40},
  breaklines=true,
  breakatwhitespace=true,
  showstringspaces=false,
  tabsize=4,
  xleftmargin=1.5em,
  framexleftmargin=1.5em,
  morekeywords={dataclass,frozen,field,Enum,IntEnum,auto,Any,
                Mapping,Sequence,Optional,match,case},
  literate={→}{{$\to$}}1 {←}{{$\leftarrow$}}1
           {≤}{{$\leq$}}1 {≥}{{$\geq$}}1 {≠}{{$\neq$}}1
           {∈}{{$\in$}}1 {∀}{{$\forall$}}1 {∃}{{$\exists$}}1
           {⊕}{{$\oplus$}}1 {⊖}{{$\ominus$}}1,
  escapeinside={(*@}{@*)},
}
\lstset{style=jugeo-python}

\lstdefinestyle{jugeo-output}{
  basicstyle=\ttfamily\small,
  backgroundcolor=\color{jugeoBg},
  frame=single,
  framerule=0.4pt,
  rulecolor=\color{jugeoGray!40},
  breaklines=true,
  showstringspaces=false,
  xleftmargin=1.5em,
  framexleftmargin=0.5em,
  numbers=none,
}

% ── JuGeo-specific macros ──────────────────────────────────────────────

% Judgment 8-tuple
\newcommand{\J}{\mathcal{J}}
\newcommand{\Jud}[8]{(#1,\,#2,\,#3,\,#4,\,#5,\,#6,\,#7,\,#8)}
\newcommand{\JudShort}[1]{\J_{#1}}

% Components
\newcommand{\coord}{\mathsf{c}}            % coordinate
\newcommand{\prop}{\varphi}                 % proposition
\newcommand{\carrier}{\mathcal{A}}          % carrier type
\newcommand{\evid}{\mathcal{E}}             % evidence bundle
\newcommand{\oblig}{\mathcal{O}}            % obligations
\newcommand{\obstr}{\mathcal{B}}            % obstructions
\newcommand{\trust}{\mathcal{T}}            % trust annotation
\newcommand{\prov}{\Pi}                     % provenance

% Trust algebra
\newcommand{\Talg}{\mathfrak{T}}            % trust ordered algebra
\newcommand{\Eadm}{\mathcal{E}_{\mathrm{adm}}}
\newcommand{\tleq}{\preceq}                 % trust ordering
\newcommand{\tjoin}{\oplus}                 % conservative join
\newcommand{\tatten}{\ominus}               % attenuation
\newcommand{\tpromote}{\uparrow_\pi}        % promotion
\newcommand{\tdemote}{\downarrow_\chi}       % demotion

% Trust tiers
\newcommand{\Tcontra}{\ensuremath{\mathsf{CONTRADICTED}}}
\newcommand{\Tunver}{\ensuremath{\mathsf{UNVERIFIED}}}
\newcommand{\Tcopilot}{\ensuremath{\mathsf{COPILOT}}}
\newcommand{\Toracle}{\ensuremath{\mathsf{ORACLE}}}
\newcommand{\Truntime}{\ensuremath{\mathsf{RUNTIME}}}
\newcommand{\Tsolver}{\ensuremath{\mathsf{SOLVER}}}
\newcommand{\Tproof}{\ensuremath{\mathsf{PROOF}}}

% Site and geometry
\newcommand{\Site}{\mathcal{S}}             % semantic site
\newcommand{\Cov}{\mathrm{Cov}}            % covering family
\newcommand{\Ob}{\mathrm{Ob}}              % objects
\newcommand{\Mor}{\mathrm{Mor}}            % morphisms
\newcommand{\res}{\mathrm{res}}            % restriction
\newcommand{\Cech}{\check{C}}              % Čech complex
\newcommand{\Hcoh}[1]{H^{#1}}             % cohomology group

% Descent
\newcommand{\descent}{\mathrm{desc}}
\newcommand{\glue}{\mathrm{glue}}
\newcommand{\overlap}[2]{#1 \cap #2}

% Semantic moves
\newcommand{\VERIFY}{\textsc{Verify}}
\newcommand{\CONSTRUCT}{\textsc{Construct}}
\newcommand{\NORMALIZE}{\textsc{Normalize}}
\newcommand{\GLUE}{\textsc{Glue}}
\newcommand{\RESTRICT}{\textsc{Restrict}}
\newcommand{\TRANSPORT}{\textsc{Transport}}
\newcommand{\REPAIR}{\textsc{Repair}}
\newcommand{\PROMOTE}{\textsc{Promote}}
\newcommand{\CHALLENGE}{\textsc{Challenge}}
\newcommand{\ESCALATE}{\textsc{Escalate}}

% Fragments
\newcommand{\QFLIA}{\texttt{QF\_LIA}}
\newcommand{\QFLRA}{\texttt{QF\_LRA}}
\newcommand{\QFBV}{\texttt{QF\_BV}}
\newcommand{\QFUF}{\texttt{QF\_UF}}

% Misc
\newcommand{\Python}{\textsc{Python}}
\newcommand{\JuGeo}{\textsc{JuGeo}}
\newcommand{\LEAN}{\textsc{Lean}\,4}
\newcommand{\Fstar}{F$^{\star}$}
\newcommand{\Coq}{\textsc{Coq}}
\newcommand{\Dafny}{\textsc{Dafny}}

% Common shortcuts
\newcommand{\ie}{i.e.\@\xspace}
\newcommand{\eg}{e.g.\@\xspace}
\newcommand{\cf}{cf.\@\xspace}
\newcommand{\wrt}{w.r.t.\@\xspace}

% Evaluation shortcuts
\newcommand{\numcases}{300}
\newcommand{\numspec}{100}
\newcommand{\numeq}{100}
\newcommand{\numbug}{100}
\newcommand{\medtime}{6.5\,ms}
\newcommand{\totaltime}{1.949\,s}


% ── Packages ────────────────────────────────────────────────────────────
\usepackage{amsmath,amssymb,amsthm,mathtools}
\usepackage{tikz-cd}
\usepackage{listings}
\usepackage{xcolor}
\usepackage{xspace}
\usepackage{booktabs}
\usepackage{multirow}
\usepackage{enumitem}
\usepackage[expansion=false]{microtype}
\usepackage{hyperref}
\usepackage{cleveref}
% \usepackage{thmtools}  % disabled: not available in this TeX installation
\usepackage{float}
\usepackage{caption}
\usepackage{subcaption}
\usepackage{tabularx}
\usepackage{stmaryrd}       % \llbracket, \rrbracket
\usepackage{bbm}            % \mathbbm{1}

% ── Page geometry ───────────────────────────────────────────────────────
\usepackage[margin=1in]{geometry}

% ── Theorem environments ────────────────────────────────────────────────
\theoremstyle{plain}
\newtheorem{theorem}{Theorem}[section]
\newtheorem{lemma}[theorem]{Lemma}
\newtheorem{proposition}[theorem]{Proposition}
\newtheorem{corollary}[theorem]{Corollary}

\theoremstyle{definition}
\newtheorem{definition}[theorem]{Definition}
\newtheorem{example}[theorem]{Example}
\newtheorem{construction}[theorem]{Construction}

\theoremstyle{remark}
\newtheorem{remark}[theorem]{Remark}
\newtheorem{notation}[theorem]{Notation}

% ── Python listings style ──────────────────────────────────────────────
\definecolor{jugeoBlue}{HTML}{2563EB}
\definecolor{jugeoGreen}{HTML}{059669}
\definecolor{jugeoRed}{HTML}{DC2626}
\definecolor{jugeoPurple}{HTML}{7C3AED}
\definecolor{jugeoGray}{HTML}{6B7280}
\definecolor{jugeoBg}{HTML}{F8FAFC}

\lstdefinestyle{jugeo-python}{
  language=Python,
  basicstyle=\ttfamily\small,
  keywordstyle=\color{jugeoBlue}\bfseries,
  stringstyle=\color{jugeoGreen},
  commentstyle=\color{jugeoGray}\itshape,
  numberstyle=\tiny\color{jugeoGray},
  backgroundcolor=\color{jugeoBg},
  numbers=left,
  numbersep=6pt,
  frame=single,
  framerule=0.4pt,
  rulecolor=\color{jugeoGray!40},
  breaklines=true,
  breakatwhitespace=true,
  showstringspaces=false,
  tabsize=4,
  xleftmargin=1.5em,
  framexleftmargin=1.5em,
  morekeywords={dataclass,frozen,field,Enum,IntEnum,auto,Any,
                Mapping,Sequence,Optional,match,case},
  literate={→}{{$\to$}}1 {←}{{$\leftarrow$}}1
           {≤}{{$\leq$}}1 {≥}{{$\geq$}}1 {≠}{{$\neq$}}1
           {∈}{{$\in$}}1 {∀}{{$\forall$}}1 {∃}{{$\exists$}}1
           {⊕}{{$\oplus$}}1 {⊖}{{$\ominus$}}1,
  escapeinside={(*@}{@*)},
}
\lstset{style=jugeo-python}

\lstdefinestyle{jugeo-output}{
  basicstyle=\ttfamily\small,
  backgroundcolor=\color{jugeoBg},
  frame=single,
  framerule=0.4pt,
  rulecolor=\color{jugeoGray!40},
  breaklines=true,
  showstringspaces=false,
  xleftmargin=1.5em,
  framexleftmargin=0.5em,
  numbers=none,
}

% ── JuGeo-specific macros ──────────────────────────────────────────────

% Judgment 8-tuple
\newcommand{\J}{\mathcal{J}}
\newcommand{\Jud}[8]{(#1,\,#2,\,#3,\,#4,\,#5,\,#6,\,#7,\,#8)}
\newcommand{\JudShort}[1]{\J_{#1}}

% Components
\newcommand{\coord}{\mathsf{c}}            % coordinate
\newcommand{\prop}{\varphi}                 % proposition
\newcommand{\carrier}{\mathcal{A}}          % carrier type
\newcommand{\evid}{\mathcal{E}}             % evidence bundle
\newcommand{\oblig}{\mathcal{O}}            % obligations
\newcommand{\obstr}{\mathcal{B}}            % obstructions
\newcommand{\trust}{\mathcal{T}}            % trust annotation
\newcommand{\prov}{\Pi}                     % provenance

% Trust algebra
\newcommand{\Talg}{\mathfrak{T}}            % trust ordered algebra
\newcommand{\Eadm}{\mathcal{E}_{\mathrm{adm}}}
\newcommand{\tleq}{\preceq}                 % trust ordering
\newcommand{\tjoin}{\oplus}                 % conservative join
\newcommand{\tatten}{\ominus}               % attenuation
\newcommand{\tpromote}{\uparrow_\pi}        % promotion
\newcommand{\tdemote}{\downarrow_\chi}       % demotion

% Trust tiers
\newcommand{\Tcontra}{\ensuremath{\mathsf{CONTRADICTED}}}
\newcommand{\Tunver}{\ensuremath{\mathsf{UNVERIFIED}}}
\newcommand{\Tcopilot}{\ensuremath{\mathsf{COPILOT}}}
\newcommand{\Toracle}{\ensuremath{\mathsf{ORACLE}}}
\newcommand{\Truntime}{\ensuremath{\mathsf{RUNTIME}}}
\newcommand{\Tsolver}{\ensuremath{\mathsf{SOLVER}}}
\newcommand{\Tproof}{\ensuremath{\mathsf{PROOF}}}

% Site and geometry
\newcommand{\Site}{\mathcal{S}}             % semantic site
\newcommand{\Cov}{\mathrm{Cov}}            % covering family
\newcommand{\Ob}{\mathrm{Ob}}              % objects
\newcommand{\Mor}{\mathrm{Mor}}            % morphisms
\newcommand{\res}{\mathrm{res}}            % restriction
\newcommand{\Cech}{\check{C}}              % Čech complex
\newcommand{\Hcoh}[1]{H^{#1}}             % cohomology group

% Descent
\newcommand{\descent}{\mathrm{desc}}
\newcommand{\glue}{\mathrm{glue}}
\newcommand{\overlap}[2]{#1 \cap #2}

% Semantic moves
\newcommand{\VERIFY}{\textsc{Verify}}
\newcommand{\CONSTRUCT}{\textsc{Construct}}
\newcommand{\NORMALIZE}{\textsc{Normalize}}
\newcommand{\GLUE}{\textsc{Glue}}
\newcommand{\RESTRICT}{\textsc{Restrict}}
\newcommand{\TRANSPORT}{\textsc{Transport}}
\newcommand{\REPAIR}{\textsc{Repair}}
\newcommand{\PROMOTE}{\textsc{Promote}}
\newcommand{\CHALLENGE}{\textsc{Challenge}}
\newcommand{\ESCALATE}{\textsc{Escalate}}

% Fragments
\newcommand{\QFLIA}{\texttt{QF\_LIA}}
\newcommand{\QFLRA}{\texttt{QF\_LRA}}
\newcommand{\QFBV}{\texttt{QF\_BV}}
\newcommand{\QFUF}{\texttt{QF\_UF}}

% Misc
\newcommand{\Python}{\textsc{Python}}
\newcommand{\JuGeo}{\textsc{JuGeo}}
\newcommand{\LEAN}{\textsc{Lean}\,4}
\newcommand{\Fstar}{F$^{\star}$}
\newcommand{\Coq}{\textsc{Coq}}
\newcommand{\Dafny}{\textsc{Dafny}}

% Common shortcuts
\newcommand{\ie}{i.e.\@\xspace}
\newcommand{\eg}{e.g.\@\xspace}
\newcommand{\cf}{cf.\@\xspace}
\newcommand{\wrt}{w.r.t.\@\xspace}

% Evaluation shortcuts
\newcommand{\numcases}{300}
\newcommand{\numspec}{100}
\newcommand{\numeq}{100}
\newcommand{\numbug}{100}
\newcommand{\medtime}{6.5\,ms}
\newcommand{\totaltime}{1.949\,s}


% ── Packages ────────────────────────────────────────────────────────────
\usepackage{amsmath,amssymb,amsthm,mathtools}
\usepackage{tikz-cd}
\usepackage{listings}
\usepackage{xcolor}
\usepackage{xspace}
\usepackage{booktabs}
\usepackage{multirow}
\usepackage{enumitem}
\usepackage[expansion=false]{microtype}
\usepackage{hyperref}
\usepackage{cleveref}
% \usepackage{thmtools}  % disabled: not available in this TeX installation
\usepackage{float}
\usepackage{caption}
\usepackage{subcaption}
\usepackage{tabularx}
\usepackage{stmaryrd}       % \llbracket, \rrbracket
\usepackage{bbm}            % \mathbbm{1}

% ── Page geometry ───────────────────────────────────────────────────────
\usepackage[margin=1in]{geometry}

% ── Theorem environments ────────────────────────────────────────────────
\theoremstyle{plain}
\newtheorem{theorem}{Theorem}[section]
\newtheorem{lemma}[theorem]{Lemma}
\newtheorem{proposition}[theorem]{Proposition}
\newtheorem{corollary}[theorem]{Corollary}

\theoremstyle{definition}
\newtheorem{definition}[theorem]{Definition}
\newtheorem{example}[theorem]{Example}
\newtheorem{construction}[theorem]{Construction}

\theoremstyle{remark}
\newtheorem{remark}[theorem]{Remark}
\newtheorem{notation}[theorem]{Notation}

% ── Python listings style ──────────────────────────────────────────────
\definecolor{jugeoBlue}{HTML}{2563EB}
\definecolor{jugeoGreen}{HTML}{059669}
\definecolor{jugeoRed}{HTML}{DC2626}
\definecolor{jugeoPurple}{HTML}{7C3AED}
\definecolor{jugeoGray}{HTML}{6B7280}
\definecolor{jugeoBg}{HTML}{F8FAFC}

\lstdefinestyle{jugeo-python}{
  language=Python,
  basicstyle=\ttfamily\small,
  keywordstyle=\color{jugeoBlue}\bfseries,
  stringstyle=\color{jugeoGreen},
  commentstyle=\color{jugeoGray}\itshape,
  numberstyle=\tiny\color{jugeoGray},
  backgroundcolor=\color{jugeoBg},
  numbers=left,
  numbersep=6pt,
  frame=single,
  framerule=0.4pt,
  rulecolor=\color{jugeoGray!40},
  breaklines=true,
  breakatwhitespace=true,
  showstringspaces=false,
  tabsize=4,
  xleftmargin=1.5em,
  framexleftmargin=1.5em,
  morekeywords={dataclass,frozen,field,Enum,IntEnum,auto,Any,
                Mapping,Sequence,Optional,match,case},
  literate={→}{{$\to$}}1 {←}{{$\leftarrow$}}1
           {≤}{{$\leq$}}1 {≥}{{$\geq$}}1 {≠}{{$\neq$}}1
           {∈}{{$\in$}}1 {∀}{{$\forall$}}1 {∃}{{$\exists$}}1
           {⊕}{{$\oplus$}}1 {⊖}{{$\ominus$}}1,
  escapeinside={(*@}{@*)},
}
\lstset{style=jugeo-python}

\lstdefinestyle{jugeo-output}{
  basicstyle=\ttfamily\small,
  backgroundcolor=\color{jugeoBg},
  frame=single,
  framerule=0.4pt,
  rulecolor=\color{jugeoGray!40},
  breaklines=true,
  showstringspaces=false,
  xleftmargin=1.5em,
  framexleftmargin=0.5em,
  numbers=none,
}

% ── JuGeo-specific macros ──────────────────────────────────────────────

% Judgment 8-tuple
\newcommand{\J}{\mathcal{J}}
\newcommand{\Jud}[8]{(#1,\,#2,\,#3,\,#4,\,#5,\,#6,\,#7,\,#8)}
\newcommand{\JudShort}[1]{\J_{#1}}

% Components
\newcommand{\coord}{\mathsf{c}}            % coordinate
\newcommand{\prop}{\varphi}                 % proposition
\newcommand{\carrier}{\mathcal{A}}          % carrier type
\newcommand{\evid}{\mathcal{E}}             % evidence bundle
\newcommand{\oblig}{\mathcal{O}}            % obligations
\newcommand{\obstr}{\mathcal{B}}            % obstructions
\newcommand{\trust}{\mathcal{T}}            % trust annotation
\newcommand{\prov}{\Pi}                     % provenance

% Trust algebra
\newcommand{\Talg}{\mathfrak{T}}            % trust ordered algebra
\newcommand{\Eadm}{\mathcal{E}_{\mathrm{adm}}}
\newcommand{\tleq}{\preceq}                 % trust ordering
\newcommand{\tjoin}{\oplus}                 % conservative join
\newcommand{\tatten}{\ominus}               % attenuation
\newcommand{\tpromote}{\uparrow_\pi}        % promotion
\newcommand{\tdemote}{\downarrow_\chi}       % demotion

% Trust tiers
\newcommand{\Tcontra}{\ensuremath{\mathsf{CONTRADICTED}}}
\newcommand{\Tunver}{\ensuremath{\mathsf{UNVERIFIED}}}
\newcommand{\Tcopilot}{\ensuremath{\mathsf{COPILOT}}}
\newcommand{\Toracle}{\ensuremath{\mathsf{ORACLE}}}
\newcommand{\Truntime}{\ensuremath{\mathsf{RUNTIME}}}
\newcommand{\Tsolver}{\ensuremath{\mathsf{SOLVER}}}
\newcommand{\Tproof}{\ensuremath{\mathsf{PROOF}}}

% Site and geometry
\newcommand{\Site}{\mathcal{S}}             % semantic site
\newcommand{\Cov}{\mathrm{Cov}}            % covering family
\newcommand{\Ob}{\mathrm{Ob}}              % objects
\newcommand{\Mor}{\mathrm{Mor}}            % morphisms
\newcommand{\res}{\mathrm{res}}            % restriction
\newcommand{\Cech}{\check{C}}              % Čech complex
\newcommand{\Hcoh}[1]{H^{#1}}             % cohomology group

% Descent
\newcommand{\descent}{\mathrm{desc}}
\newcommand{\glue}{\mathrm{glue}}
\newcommand{\overlap}[2]{#1 \cap #2}

% Semantic moves
\newcommand{\VERIFY}{\textsc{Verify}}
\newcommand{\CONSTRUCT}{\textsc{Construct}}
\newcommand{\NORMALIZE}{\textsc{Normalize}}
\newcommand{\GLUE}{\textsc{Glue}}
\newcommand{\RESTRICT}{\textsc{Restrict}}
\newcommand{\TRANSPORT}{\textsc{Transport}}
\newcommand{\REPAIR}{\textsc{Repair}}
\newcommand{\PROMOTE}{\textsc{Promote}}
\newcommand{\CHALLENGE}{\textsc{Challenge}}
\newcommand{\ESCALATE}{\textsc{Escalate}}

% Fragments
\newcommand{\QFLIA}{\texttt{QF\_LIA}}
\newcommand{\QFLRA}{\texttt{QF\_LRA}}
\newcommand{\QFBV}{\texttt{QF\_BV}}
\newcommand{\QFUF}{\texttt{QF\_UF}}

% Misc
\newcommand{\Python}{\textsc{Python}}
\newcommand{\JuGeo}{\textsc{JuGeo}}
\newcommand{\LEAN}{\textsc{Lean}\,4}
\newcommand{\Fstar}{F$^{\star}$}
\newcommand{\Coq}{\textsc{Coq}}
\newcommand{\Dafny}{\textsc{Dafny}}

% Common shortcuts
\newcommand{\ie}{i.e.\@\xspace}
\newcommand{\eg}{e.g.\@\xspace}
\newcommand{\cf}{cf.\@\xspace}
\newcommand{\wrt}{w.r.t.\@\xspace}

% Evaluation shortcuts
\newcommand{\numcases}{300}
\newcommand{\numspec}{100}
\newcommand{\numeq}{100}
\newcommand{\numbug}{100}
\newcommand{\medtime}{6.5\,ms}
\newcommand{\totaltime}{1.949\,s}


\usepackage{xspace}          % safe spacing after macros

% ── Paper-specific macros ─────────────────────────────────────────────
\newcommand{\treaty}{\mathcal{T}}
\newcommand{\patch}{\mathcal{P}}
\newcommand{\iface}{\mathcal{I}}
\newcommand{\exports}{\mathrm{exp}}
\newcommand{\imports}{\mathrm{imp}}
\newcommand{\shared}{\mathrm{sh}}
\newcommand{\cscore}[1][]{c_{\mathrm{score}#1}}
\newcommand{\cthresh}{\theta}
\newcommand{\negotiate}{\textsc{Negotiate}}
\newcommand{\synth}{\textsc{Synthesize}}
\newcommand{\detect}{\textsc{Detect}}
\newcommand{\resolve}{\textsc{Resolve}}
\newcommand{\HC}{\mathrm{HC}}
\newcommand{\Nerve}{\mathrm{N}}
\newcommand{\cosk}{\mathrm{cosk}}

\title{\textbf{Automated Interface Reconciliation for Modular Verification}}
\author{%
  Halley Young\\
  \textit{Microsoft Research}\\
  \texttt{halleyyoung@microsoft.com}
}
\date{\today}

\begin{document}

\maketitle

\begin{abstract}
When a program is decomposed into overlapping patches for modular
verification, the patches must agree on shared boundaries---but different
patches may export conflicting interfaces.  We introduce \emph{treaty
synthesis}, a dynamic negotiation protocol that reconciles patch interfaces
within the \emph{hypercover family} framework of Judgment Geometry.  A
treaty is a formal witness that the \v{C}ech $1$-cocycle on patch
overlaps vanishes, enabling descent from local verifications to a global
correctness certificate.  Our synthesis algorithm iterates through
conflict detection, scoring, and resolution---applying strategies such as
common-supertype widening, range intersection, protocol adaptation, and
ownership brokering---with guaranteed termination in at most $10$
negotiation rounds via a strictly decreasing obstruction norm.  We
evaluate treaty synthesis on six conflict scenarios drawn from
real-world interface reconciliation patterns
(\textsf{type\_mismatch}, \textsf{range\_conflict},
\textsf{nullable\_disagreement}, \textsf{protocol\_version},
\textsf{encoding\_mismatch}, \textsf{resource\_lifetime}),
detecting $18$ conflicts and resolving $10$ of them ($55.6\%$
resolution rate).  The negotiation protocol converges in an average of
$8.83$ rounds; four of six scenarios achieve full convergence
(obstruction norm below $10^{-6}$), while two hit the $10$-round cap.
The obstruction norm decays geometrically at a factor of approximately
$0.2$ per round, and all scenarios reach \textsf{PROOF\_BACKED} trust.
Compared to ML module signatures~\cite{milner1997,leroy1994}, Haskell
type-class instances~\cite{wadler1989}, and \Fstar{}
interfaces~\cite{swamy2016}, treaty synthesis is the first mechanism to
support dynamic negotiation, conflict scoring, trust propagation, and
automatic repair on failure.
\end{abstract}

% ======================================================================
\section{Introduction}\label{sec:intro}
% ======================================================================

Modular verification decomposes a large program into manageable pieces,
verifies each piece independently, and then assembles the results into a
global correctness argument.  The critical step---often the hardest---is
ensuring that the independently verified pieces \emph{agree} on their
shared boundaries.  In type theory, this is the domain of module systems:
ML signatures~\cite{milner1997,leroy1994} prescribe interfaces that modules
must satisfy, and the type-checker enforces compatibility statically.  In
verification, \Fstar{} interfaces~\cite{swamy2016} and \Dafny{}
refinement types~\cite{leino2010} play an analogous role.

All of these approaches are \emph{static}: the interface must be fixed
before verification begins, and any mismatch is a hard error.  This is
adequate when the decomposition is controlled by a single designer, but
it breaks down in the setting we care about: \emph{automatic}
decomposition of programs into overlapping patches by a tool that may not
know, a priori, which interfaces the patches will export.

\paragraph{The treaty synthesis problem.}
Consider two patches $\patch_A$ and $\patch_B$ that overlap on a shared
region $U_{AB}$.  Patch $\patch_A$ exports a symbol \texttt{solve} with
type $\mathtt{int} \to \mathtt{int}$; patch $\patch_B$ exports the same
symbol with type $\mathtt{list} \to \mathtt{int}$.  A static module
system would simply reject the composition.  But the conflict may be
\emph{resolvable}: perhaps the patches agree on the subset of inputs
that actually flow through $U_{AB}$, or perhaps one patch's type is a
subtype of the other.  Treaty synthesis detects such conflicts, scores
their severity, selects a resolution strategy, and iterates until the
conflict score drops below a threshold---or declares the composition
irreconcilable and reports a cohomological obstruction.

\paragraph{Contributions.}
\begin{enumerate}[label=(\roman*),nosep]
\item We define \emph{hypercover families} (\Cref{sec:hypercover}), a
  higher-order decomposition framework that supports module $\to$ class
  $\to$ method $\to$ branch decomposition.
\item We present the \emph{treaty synthesis} algorithm
  (\Cref{sec:synthesis}) implemented in the \JuGeo{} system, with
  formal conflict detection (\Cref{sec:conflict}), four resolution
  strategies (\Cref{sec:strategies}), and a negotiation loop with
  guaranteed termination (\Cref{sec:negotiation}).
\item We connect treaty resolution to \v{C}ech cohomology: a resolved
  treaty witnesses the vanishing of $\Hcoh{1}$, and treaty failure
  produces an obstruction class (\Cref{sec:integration}).
\item We evaluate on six conflict scenarios from real-world interface
  reconciliation patterns, achieving a $66.7\%$ full convergence rate
  ($4$ of $6$ scenarios) with geometric obstruction-norm decay, and
  compare to four state-of-the-art interface mechanisms
  (\Cref{sec:evaluation}).
\end{enumerate}

\paragraph{Outline.}
\Cref{sec:hypercover} introduces hypercover families.
\Cref{sec:synthesis} presents the synthesis algorithm.
\Cref{sec:conflict} details conflict detection and scoring.
\Cref{sec:strategies} describes the four resolution strategies.
\Cref{sec:negotiation} proves termination of the negotiation loop.
\Cref{sec:integration} connects treaties to descent and trust.
\Cref{sec:evaluation} presents the evaluation.
\Cref{sec:related} surveys related work.
\Cref{sec:conclusion} concludes.

\paragraph{Motivating scenario.}
To ground the discussion, consider a program $P$ that implements a
sorting algorithm.  A verification tool decomposes $P$ into three
patches: $\patch_1$ covers the partitioning logic, $\patch_2$ covers
the recursive calls, and $\patch_3$ covers the base case.  The overlaps
$\patch_1 \cap \patch_2$ and $\patch_2 \cap \patch_3$ involve shared
variables (the pivot element, the array bounds).  If $\patch_1$ assumes
the pivot is an \texttt{int} but $\patch_2$ treats it as a
\texttt{float} (due to a polymorphic comparison), we have an interface
contradiction.  Treaty synthesis detects this, scores it (say $0.4$,
moderate severity since \texttt{int}${}\hookrightarrow{}$\texttt{float}
is a safe coercion), and resolves it via \textsc{Merge} (widening to
\texttt{float}).  The following diagram shows the three-patch scenario:

\begin{center}
\begin{tikzcd}[column sep=2.5em, row sep=1.8em]
  \patch_1 \arrow[dr, "\treaty_{12}"']
  & & \patch_3 \arrow[dl, "\treaty_{23}"] \\
  & \patch_2 \arrow[ul, bend left, dashed, "\shared_{12}"]
    \arrow[ur, bend right, dashed, "\shared_{23}"'] &
\end{tikzcd}
\end{center}

\noindent The dashed arrows indicate shared interfaces; the solid arrows
indicate treaties.  The treaty $\treaty_{12}$ resolves the
\texttt{int}/\texttt{float} conflict; $\treaty_{23}$ is trivially
resolved (no conflicts).  The global correctness argument assembles via
descent over the nerve of the cover $\{\patch_1, \patch_2, \patch_3\}$.


% ======================================================================
\section{Hypercover Families}\label{sec:hypercover}
% ======================================================================

A \emph{cover} of a program $P$ is a collection of patches
$\{U_i\}_{i \in I}$ such that $P = \bigcup_{i \in I} U_i$.  A
\emph{hypercover} is a cover of a cover: a two-level decomposition in
which each patch $U_i$ is itself covered by sub-patches
$\{V_{ij}\}_{j \in J_i}$.  More generally:

\begin{definition}[Hypercover family]\label{def:hypercover}
Let $\Site = (\mathsf{Prog}, \Cov)$ be the semantic site of programs
with covering families.  A \emph{hypercover family} of level $n$ over a
program $P$ is an augmented simplicial object
\[
  \HC_\bullet \;:\;
  \HC_n \rightrightarrows \cdots \rightrightarrows \HC_1
  \rightrightarrows \HC_0 \to P
\]
such that at each level $k$, the object $\HC_k$ is a covering family for
the $(k{-}1)$-fold overlaps of $\HC_{k-1}$, and the natural map
$\HC_k \to (\cosk_{k-1}\,\HC)_k$ is a cover.
\end{definition}

In practice, we use hypercovers of level $\leq 3$:

\begin{center}
\begin{tikzcd}[row sep=2.2em, column sep=1.8em]
  V_{11} \arrow[d] & V_{12} \arrow[d] & V_{21} \arrow[d] & V_{22} \arrow[d] \\
  U_1 \arrow[dr] & U_2 \arrow[d] & U_3 \arrow[dl] & U_4 \arrow[dll] \\
  & C_1 \arrow[d] & C_2 \arrow[dl] & \\
  & P & &
\end{tikzcd}
\end{center}

\noindent The nodes at each level form a covering family: the $V_{ij}$
cover the methods $U_k$, the methods cover the classes $C_l$, and the
classes cover the module $P$.  Each arrow represents an inclusion
(restriction) map.  The key feature of a hypercover, as opposed to a
mere iterated cover, is the \emph{descent condition}: the covering at
each level must be sufficiently fine to detect all the information
present at the level below.

\paragraph{Why hypercovers?}
A simple cover suffices for small programs, but realistic software has
hierarchical structure: modules contain classes, classes contain methods,
methods contain branches.  A hypercover family captures this hierarchy
and allows verification to proceed \emph{level by level}.  At each
level, the patches may overlap, and overlapping patches must agree on
their shared interfaces.  This is where treaty synthesis enters.

\begin{example}[Three-level hypercover]\label{ex:hypercover}
Consider a Python module \texttt{math\_utils.py} containing two classes,
each with two methods.  The level-$0$ cover consists of the two classes:
$\mathcal{U} = \{C_1, C_2\}$.  The level-$1$ cover refines each class
into its methods: $C_1 = U_1 \cup U_2$ and $C_2 = U_3 \cup U_4$.
The level-$2$ cover refines each method into its branches:
$U_1 = V_{11} \cup V_{12}$ (the \texttt{if} and \texttt{else} branches
of the first method).

The overlap $C_1 \cap C_2$ consists of symbols shared between the two
classes---utility functions, shared state, or inherited methods.  At the
method level, $U_1 \cap U_2$ consists of local variables or helper calls
shared between two methods of the same class.  Treaty synthesis operates
independently at each level, starting from the finest (branches) and
propagating upward.
\end{example}

\paragraph{Coskeletal condition.}
The defining property of a hypercover---that $\HC_k \to (\cosk_{k-1}\,\HC)_k$
is a cover---ensures that the refinement at level $k$ is ``fine enough''
relative to the lower levels.  Concretely, this means that every
$(k{-}1)$-fold overlap that arises from level $k{-}1$ is already covered
by patches at level $k$.  Without this condition, we could have
overlaps at level $k{-}1$ that are invisible at level $k$, and treaty
synthesis would miss potential conflicts.

\begin{definition}[Nerve of a hypercover]\label{def:nerve}
The \emph{nerve} $\Nerve(\HC_k)$ of a level-$k$ hypercover is the
simplicial complex whose $p$-simplices are $(p{+}1)$-fold overlaps
$U_{i_0} \cap \cdots \cap U_{i_p} \neq \emptyset$.  The \v{C}ech
complex $\Cech^\bullet(\HC_k, \mathcal{F})$ for a presheaf $\mathcal{F}$
is the cochain complex on $\Nerve(\HC_k)$.
\end{definition}

The connection to treaty synthesis is:

\begin{proposition}\label{prop:treaty-cocycle}
A treaty on the overlap $U_i \cap U_j$ is a section of the interface
presheaf $\iface$ over $U_i \cap U_j$.  The collection of all treaties
is a \v{C}ech $1$-cochain $\treaty \in \Cech^1(\HC_0, \iface)$.  The
treaties are \emph{consistent} if and only if $\treaty$ is a cocycle,
\ie{} $\delta^1 \treaty = 0$.
\end{proposition}

The following diagram summarizes the relationship between hypercovers,
treaties, and cohomology:

\begin{equation}\label{eq:cech-sequence}
\begin{tikzcd}[column sep=2.8em]
  0 \arrow[r]
  & \Cech^0(\HC, \iface) \arrow[r, "\delta^0"]
  & \Cech^1(\HC, \iface) \arrow[r, "\delta^1"]
  & \Cech^2(\HC, \iface) \\[-0.5em]
  & \text{\small (local interfaces)}
  & \text{\small (treaties)}
  & \text{\small (triple conflicts)}
\end{tikzcd}
\end{equation}

A resolved treaty lives in $\ker \delta^1$; if it is also in
$\mathrm{im}\,\delta^0$, the treaty is \emph{trivial} (the patches
already agreed).  The obstruction to global agreement is the cohomology
class $[\treaty] \in \Hcoh{1}(\HC, \iface)$.


% ======================================================================
\section{Treaty Synthesis}\label{sec:synthesis}
% ======================================================================

We now describe the treaty synthesis algorithm as implemented in the
\JuGeo{} system.  The implementation resides in
\texttt{src/jugeo/generation/hypercover\_treaties/algorithms.py}.

\subsection{Data Model}

A \emph{patch interface} is a record describing what a patch exports and
imports:

\begin{lstlisting}[caption={Patch interface descriptor and resolution strategies.},
  label=lst:interface]
@dataclass(frozen=True)
class PatchInterface:
    patch_id: str
    exports: Mapping[str, str]   # symbol → type signature
    imports: Mapping[str, str]   # symbol → expected type
    version: str                 # semantic version
    trust: TrustTier             # current trust level

class ResolutionStrategy(str, Enum):
    PREFER_LEFT = "PREFER_LEFT"     # adopt treaty A
    PREFER_RIGHT = "PREFER_RIGHT"   # adopt treaty B
    MERGE = "MERGE"                 # union of exports, meet of trust
    SPLIT = "SPLIT"                 # quarantine conflicting symbols

MAX_NEGOTIATION_ROUNDS = 10
CONFLICT_SCORE_THRESHOLD = 0.3
\end{lstlisting}

A \emph{treaty} formalizes the agreement between two patches:

\begin{definition}[Treaty]\label{def:treaty}
A treaty between patches $\patch_A$ and $\patch_B$ is a $5$-tuple
\[
  \treaty = (\iface_A,\; \iface_B,\; S,\; E,\; R)
\]
where $\iface_A, \iface_B$ are the patch interfaces, $S$ is the shared
interface (symbols both patches reference), $E$ is the merged export
map, and $R$ is the resolution record.
\end{definition}

\subsection{The \texttt{TreatySynthesizer}}

The synthesizer takes two patch interfaces and produces a treaty:

\begin{lstlisting}[caption={Treaty synthesis: shared interface and export map computation.},
  label=lst:synthesizer]
class TreatySynthesizer:
    def synthesize(self, iface_a: PatchInterface,
                   iface_b: PatchInterface) -> Treaty:
        # Shared interface: symbols that cross the boundary
        shared = {}
        for sym in iface_a.exports:
            if sym in iface_b.imports:
                shared[sym] = iface_a.exports[sym]
        for sym in iface_b.exports:
            if sym in iface_a.imports:
                shared[sym] = iface_b.exports[sym]

        # Merged export map
        export_map = {**iface_a.exports, **iface_b.exports}

        # Treaty hash for identity tracking
        treaty_hash = hashlib.sha256(
            json.dumps(shared, sort_keys=True).encode()
        ).hexdigest()[:16]

        return Treaty(
            iface_a=iface_a, iface_b=iface_b,
            shared=shared, export_map=export_map,
            treaty_hash=treaty_hash,
        )
\end{lstlisting}

\begin{remark}
The shared interface $S$ is computed as
$S = (\exports_A \cap \imports_B) \cup (\exports_B \cap \imports_A)$.
This is exactly the set of symbols that flow \emph{across} the patch
boundary---the symbols relevant for the overlap $U_A \cap U_B$.
\end{remark}

The synthesis process is captured by the following diagram:

\begin{equation}\label{eq:synthesis-diagram}
\begin{tikzcd}[column sep=3em, row sep=2em]
  \iface_A \arrow[dr, "\exports_A"'] \arrow[r, dashed, "\shared"]
  & U_A \cap U_B \arrow[d, "\treaty"]
  & \iface_B \arrow[dl, "\exports_B"] \arrow[l, dashed, "\shared"'] \\
  & E
\end{tikzcd}
\end{equation}

\subsection{Formal Properties}

\begin{proposition}[Synthesis soundness]\label{prop:synth-sound}
If $\shared(\iface_A, \iface_B) = \emptyset$, then
$\treaty = (\iface_A, \iface_B, \emptyset, \exports_A \cup \exports_B, \bot)$
is trivially resolved---the patches do not interact and no negotiation
is required.
\end{proposition}

\begin{proof}
If the shared interface is empty, no symbol crosses the patch boundary.
The merged export map is the disjoint union, and the $1$-cocycle is
the zero cochain.
\end{proof}

\begin{proposition}[Synthesis completeness]\label{prop:synth-complete}
For any two patch interfaces $\iface_A, \iface_B$ over a common overlap
$U_A \cap U_B \neq \emptyset$, the synthesizer produces a treaty.  The
treaty may contain unresolved conflicts, but it is always well-defined.
\end{proposition}

\begin{proof}
The shared interface computation terminates in $O(|\exports_A| +
|\exports_B|)$ time, and the export map is always defined as the union.
Conflicts are recorded but do not prevent treaty construction.
\end{proof}


% ======================================================================
\section{Conflict Detection}\label{sec:conflict}
% ======================================================================

Not all treaties are conflict-free.  The \texttt{ConflictDetector}
identifies three types of conflicts.

\subsection{Conflict Taxonomy}

\begin{definition}[Conflict types]\label{def:conflict-types}
Let $\treaty = (\iface_A, \iface_B, S, E, R)$ be a treaty.  A
\emph{conflict} is one of:
\begin{enumerate}[label=(\alph*),nosep]
\item \textbf{Interface contradiction}: a symbol $s \in S$ such that
  $\iface_A.\exports(s) \neq \iface_B.\exports(s)$ and neither type
  is a subtype of the other.
\item \textbf{Export overlap}: a symbol $s$ such that both
  $s \in \exports_A$ and $s \in \exports_B$ but $s \notin S$---both
  patches claim to provide $s$ without either importing it.
\item \textbf{Version mismatch}: the patches declare incompatible
  dependency versions for a shared library.
\end{enumerate}
\end{definition}

Each conflict is assigned a \emph{severity score}
$\cscore \in [0, 1]$:

\begin{lstlisting}[caption={Conflict detection with severity scoring.},
  label=lst:detector]
class ConflictDetector:
    def detect(self, treaty: Treaty) -> list[Conflict]:
        conflicts = []
        for sym, type_a in treaty.iface_a.exports.items():
            if sym in treaty.iface_b.exports:
                type_b = treaty.iface_b.exports[sym]
                if type_a != type_b:
                    # Interface contradiction
                    score = self._type_distance(type_a, type_b)
                    conflicts.append(Conflict(
                        kind="INTERFACE_CONTRADICTION",
                        symbol=sym,
                        left_type=type_a, right_type=type_b,
                        score=score,
                    ))
                # else: compatible, no conflict

        # Export overlap detection
        for sym in treaty.iface_a.exports:
            if sym in treaty.iface_b.exports:
                if sym not in treaty.shared:
                    conflicts.append(Conflict(
                        kind="EXPORT_OVERLAP",
                        symbol=sym, score=0.5,
                    ))

        # Version mismatch detection
        if treaty.iface_a.version != treaty.iface_b.version:
            if not self._versions_compatible(
                treaty.iface_a.version,
                treaty.iface_b.version,
            ):
                conflicts.append(Conflict(
                    kind="VERSION_MISMATCH",
                    score=0.6,
                ))

        return conflicts
\end{lstlisting}

\subsection{Concrete Example}

Consider two patches with the following exports:

\begin{lstlisting}[caption={Conflict example: incompatible \texttt{solve} types.},
  label=lst:conflict-example]
# Patch A exports: {solve: int -> int, helper: str -> str}
# Patch B exports: {solve: list -> int, helper: str -> str}

iface_a = PatchInterface(
    patch_id="A",
    exports={"solve": "int -> int", "helper": "str -> str"},
    imports={"solve": "int -> int"},
    version="1.0.0", trust=TrustTier.VERIFIED,
)
iface_b = PatchInterface(
    patch_id="B",
    exports={"solve": "list -> int", "helper": "str -> str"},
    imports={"solve": "list -> int"},
    version="1.0.0", trust=TrustTier.VERIFIED,
)

# ConflictDetector output:
# [Conflict(kind="INTERFACE_CONTRADICTION",
#           symbol="solve",
#           left_type="int -> int",
#           right_type="list -> int",
#           score=0.8)]
# helper: compatible (no conflict)
\end{lstlisting}

The score $0.8$ reflects high severity: the input types
$\mathtt{int}$ and $\mathtt{list}$ are unrelated (neither is a subtype
of the other).  By contrast, a conflict between $\mathtt{int}$ and
$\mathtt{float}$ might score $0.3$ due to implicit coercion.

\subsection{Aggregate Conflict Score}

The aggregate conflict score for a treaty is defined as:

\begin{definition}[Aggregate conflict score]\label{def:agg-score}
Given conflicts $c_1, \ldots, c_m$ with scores
$\cscore[1], \ldots, \cscore[m]$, the aggregate score is
\[
  \cscore(\treaty) = 1 - \prod_{k=1}^{m}(1 - \cscore[k]).
\]
This is the probabilistic-or of the individual scores: the treaty is
conflict-free ($\cscore = 0$) only if every individual score is $0$.
\end{definition}

\begin{lemma}\label{lem:score-monotone}
If a resolution step eliminates conflict $c_j$ (setting $\cscore[j] = 0$)
without introducing new conflicts, then $\cscore(\treaty)$ strictly
decreases.
\end{lemma}

\begin{proof}
Removing one factor from the product strictly increases
$\prod_k (1 - \cscore[k])$, hence strictly decreases
$1 - \prod_k (1 - \cscore[k])$.
\end{proof}


% ======================================================================
\section{Resolution Strategies}\label{sec:strategies}
% ======================================================================

Once conflicts are detected, they must be resolved.  The \JuGeo{} system
supports four resolution strategies, each with precise formal semantics.

\begin{definition}[Resolution strategies]\label{def:strategies}
Let $\treaty = (\iface_A, \iface_B, S, E, R)$ be a treaty with
conflict $c$ on symbol $s$.  The strategies are:

\begin{enumerate}[label=(\roman*),nosep]
\item \textbf{\textsc{Prefer-Left}}: adopt $\iface_A$'s type for $s$.
  Formally, $E' = E[s \mapsto \iface_A.\exports(s)]$.  Trust is
  inherited from $\patch_A$.

\item \textbf{\textsc{Prefer-Right}}: adopt $\iface_B$'s type for $s$.
  Formally, $E' = E[s \mapsto \iface_B.\exports(s)]$.  Trust is
  inherited from $\patch_B$.

\item \textbf{\textsc{Merge}}: take the union of exports and the
  \emph{meet} (greatest lower bound) of trust.  The merged type for $s$
  is the least common supertype if it exists, or a tagged union
  otherwise.  Formally:
  \[
    E' = E[s \mapsto \iface_A.\exports(s) \cup \iface_B.\exports(s)],
    \quad
    \trust' = \trust_A \wedge \trust_B.
  \]

\item \textbf{\textsc{Split}}: quarantine the conflicting symbol $s$
  into a separate boundary module.  The treaty records $s$ as
  \emph{unresolved} and prevents it from flowing across the patch
  boundary.  Formally, $S' = S \setminus \{s\}$ and
  $E' = E \setminus \{s\}$.
\end{enumerate}
\end{definition}

The strategies form a lattice ordered by \emph{information content}:

\begin{center}
\begin{tikzcd}[row sep=1.5em, column sep=0.5em]
  & \textsc{Merge} \arrow[dl, dash] \arrow[dr, dash] & \\
  \textsc{Prefer-Left} \arrow[dr, dash]
  && \textsc{Prefer-Right} \arrow[dl, dash] \\
  & \textsc{Split} &
\end{tikzcd}
\end{center}

\textsc{Merge} is the richest strategy (most symbols survive) but may
lower trust.  \textsc{Split} is the most conservative (fewest symbols,
highest residual trust).  The system defaults to \textsc{Merge} and
falls back to \textsc{Split} when the conflict score exceeds $0.7$.

\begin{proposition}[Strategy soundness]\label{prop:strategy-sound}
Each strategy produces a treaty $\treaty'$ with
$\cscore(\treaty') \leq \cscore(\treaty)$.  Equality holds only for
\textsc{Split} when the quarantined symbol reappears in a transitive
dependency.
\end{proposition}

\begin{proof}
\textsc{Prefer-Left} and \textsc{Prefer-Right} eliminate the conflict on
$s$ by fiat, setting $\cscore[s] = 0$.  \textsc{Merge} resolves the
conflict if a common supertype exists; otherwise it falls back to
\textsc{Split}.  \textsc{Split} removes $s$ from $S$, so the conflict on
$s$ no longer exists.  In all cases, the aggregate score does not
increase by \Cref{lem:score-monotone}.
\end{proof}

\begin{example}[Strategy comparison]\label{ex:strategy-comparison}
Consider the conflict from \Cref{lst:conflict-example}: \texttt{solve}
has type \texttt{int -> int} in $\patch_A$ and \texttt{list -> int} in
$\patch_B$, with score $0.8$.

\begin{itemize}[nosep]
\item \textsc{Prefer-Left}: the treaty exports
  $\texttt{solve}: \texttt{int -> int}$.  Patch $B$ must adapt its
  usage.  Trust inherits from $\patch_A$.
\item \textsc{Prefer-Right}: the treaty exports
  $\texttt{solve}: \texttt{list -> int}$.  Patch $A$ must adapt.
  Trust inherits from $\patch_B$.
\item \textsc{Merge}: the system attempts to find
  $\texttt{int} \sqcup \texttt{list}$.  Since no common supertype
  exists in the type lattice, \textsc{Merge} falls back to a tagged
  union $\texttt{int | list} \to \texttt{int}$, with trust demoted to
  $\trust_A \wedge \trust_B$.
\item \textsc{Split}: \texttt{solve} is quarantined---removed from the
  shared interface.  Each patch retains its own version, and neither
  flows across the boundary.
\end{itemize}

\noindent For this conflict (score $\geq 0.7$), the heuristic selects
\textsc{Split}, which is the safest choice when the types are
fundamentally incompatible.
\end{example}


% ======================================================================
\section{The Negotiation Loop}\label{sec:negotiation}
% ======================================================================

Treaty synthesis, conflict detection, and resolution are composed into a
\emph{negotiation loop} that iterates until all conflicts are resolved
or the maximum number of rounds is reached.

\subsection{The \texttt{TreatyNegotiator}}

\begin{lstlisting}[caption={Negotiation loop: synthesize, detect, resolve, re-check.},
  label=lst:negotiator]
class TreatyNegotiator:
    def __init__(self):
        self.synthesizer = TreatySynthesizer()
        self.detector = ConflictDetector()
        self.max_rounds = MAX_NEGOTIATION_ROUNDS  # 10

    def negotiate(self, iface_a: PatchInterface,
                  iface_b: PatchInterface) -> NegotiationResult:
        treaty = self.synthesizer.synthesize(iface_a, iface_b)
        prev_score = float('inf')

        for round_num in range(self.max_rounds):
            conflicts = self.detector.detect(treaty)
            score = self._aggregate_score(conflicts)

            if score == 0.0:
                return NegotiationResult(
                    status="resolved", treaty=treaty,
                    rounds=round_num + 1,
                )

            if score < CONFLICT_SCORE_THRESHOLD:  # 0.3
                return NegotiationResult(
                    status="soft_resolved", treaty=treaty,
                    rounds=round_num + 1,
                    residual_score=score,
                )

            assert score < prev_score, "score must strictly decrease"
            prev_score = score

            # Apply resolution strategy per conflict
            for conflict in conflicts:
                strategy = self._select_strategy(conflict)
                treaty = self._apply_strategy(
                    treaty, conflict, strategy
                )

        return NegotiationResult(
            status="unresolved", treaty=treaty,
            rounds=self.max_rounds,
        )
\end{lstlisting}

\subsection{Termination Proof}

\begin{theorem}[Termination]\label{thm:termination}
The negotiation loop terminates in at most
$N = \texttt{MAX\_NEGOTIATION\_ROUNDS} = 10$ iterations.  If no external
conflicts are introduced, the conflict score strictly decreases at each
round.
\end{theorem}

\begin{proof}
We establish two claims.

\emph{Claim 1: Bounded iterations.}  The loop is bounded by
$N = 10$, enforced by the \texttt{for} loop.  This is a hard upper
bound regardless of conflict behavior.

\emph{Claim 2: Strict decrease.}  At each round, the negotiator applies
a resolution strategy to every detected conflict.  By
\Cref{prop:strategy-sound}, each strategy produces a treaty with
conflict score at most that of the previous round.  The inequality is
strict because at least one conflict is fully resolved
($\cscore[j] \to 0$) or quarantined ($s$ removed from $S$), and by
\Cref{lem:score-monotone} the aggregate score strictly decreases.

The assertion \texttt{assert score < prev\_score} serves as a runtime
check of Claim~2.  If it fires, the system has a bug in strategy
application.

Combining the claims: the loop either terminates early (score reaches
$0$ or drops below $\cthresh = 0.3$) or exhausts $N$ rounds.
\end{proof}

\begin{lemma}[Geometric Convergence of Obstruction Norm]\label{lem:geom-convergence}
Let $\|\obstr_k\|$ denote the obstruction norm at round $k$.  If the
resolution strategy at each round reduces the number of violated overlaps
by a constant fraction $\rho \in (0,1)$, then
\[
  \|\obstr_k\| \;\leq\; \|\obstr_0\| \cdot \rho^k
\]
and convergence (to norm below threshold $\epsilon$) requires at most
$\lceil \log(\|\obstr_0\|/\epsilon) / \log(1/\rho) \rceil$ rounds.
\end{lemma}

\begin{proof}
At each round, the negotiator applies a resolution strategy to every detected
conflict.  By \Cref{prop:strategy-sound}, each application reduces
the aggregate score.  The obstruction norm is defined as the sum of
conflict-weighted overlap violations:
\[
  \|\obstr_k\| = \sum_{(i,j) \in \text{violated}} w_{ij} \cdot d(\res(\sigma_i), \res(\sigma_j))
\]
where $d$ is a type-metric distance and $w_{ij}$ are overlap weights.
If the strategy reduces each term by factor $\rho$, the sum decreases by
$\rho$, giving $\|\obstr_{k+1}\| \leq \rho \cdot \|\obstr_k\|$.
Iterating, $\|\obstr_k\| \leq \|\obstr_0\| \cdot \rho^k$.
Setting $\|\obstr_0\| \cdot \rho^k < \epsilon$ and solving for $k$ gives
the bound.

Empirically, we observe $\rho \approx 0.2$ for the \textsc{Widen} strategy
and $\rho \approx 0.15$ for the \textsc{Version-Adapter} strategy
(\Cref{tab:convergence}), consistent with the geometric bound.
\end{proof}

\begin{corollary}\label{cor:complexity}
The worst-case complexity of negotiation is
$O(N \cdot |\exports_A| \cdot |\exports_B|)$, where $N = 10$, and the
factor $|\exports_A| \cdot |\exports_B|$ accounts for conflict detection.
\end{corollary}

\subsection{Strategy Selection}

The negotiator selects strategies using the following heuristic:

\begin{enumerate}[nosep]
\item If $\cscore < 0.3$: use \textsc{Merge} (soft conflict, likely
  resolvable by supertyping).
\item If $0.3 \leq \cscore < 0.7$: use \textsc{Prefer-Left} or
  \textsc{Prefer-Right} based on trust comparison
  ($\trust_A \tleq \trust_B$ implies prefer $B$).
\item If $\cscore \geq 0.7$: use \textsc{Split} (hard conflict,
  quarantine the symbol).
\end{enumerate}

This yields a monotone refinement: each round either merges soft
conflicts, resolves moderate ones by trust, or quarantines hard ones.


% ======================================================================
\section{Integration with Descent and Trust}\label{sec:integration}
% ======================================================================

Treaty synthesis does not operate in isolation.  It integrates with the
\emph{descent} framework of Judgment Geometry (which glues local
verifications into global correctness) and the \emph{trust algebra}
(which tracks the epistemic quality of each verification step).

\subsection{Treaties as Cocycle Witnesses}

Recall from \Cref{prop:treaty-cocycle} that a collection of treaties
forms a \v{C}ech $1$-cochain $\treaty \in \Cech^1(\HC_0, \iface)$.
The key insight is:

\begin{theorem}[Treaty--descent correspondence]\label{thm:treaty-descent}
A fully resolved treaty (conflict score $0$) is a witness that the
\v{C}ech $1$-cocycle $\delta^0(\sigma) = \treaty$ vanishes in
cohomology.  Equivalently, the local interface sections $\sigma_i \in
\iface(U_i)$ glue to a global section $\sigma \in \iface(P)$.
\end{theorem}

\begin{proof}[Proof sketch]
A resolved treaty means that for every overlap $U_i \cap U_j$, the
patches agree on the shared interface:
$\res_{U_i \cap U_j}(\sigma_i) = \res_{U_i \cap U_j}(\sigma_j)$.  This
is exactly the cocycle condition $\delta^0(\sigma)_{ij} = \sigma_j -
\sigma_i = 0$.  By the sheaf condition on $\iface$, the local sections
glue.
\end{proof}

\begin{theorem}[Treaty Soundness]\label{thm:treaty-sound}
Let $\{\patch_i\}_{i \in I}$ be a hypercover of program $P$, and let
$\{\treaty_{ij}\}$ be the collection of resolved treaties (obstruction
norm $< \epsilon$ for threshold $\epsilon > 0$).  Then there exists a
global interface section $\sigma \in \iface(P)$ such that
\[
  \|\res_{U_i \cap U_j}(\sigma) - \treaty_{ij}\| < \epsilon
  \quad \text{for all } (i,j).
\]
Moreover, the descent from local verifications to a global correctness
certificate is compatible with the treaty-resolved interfaces.
\end{theorem}

\begin{proof}
The collection $\{\treaty_{ij}\}$ forms a \v{C}ech $1$-cochain.
Since each treaty is resolved (its obstruction norm is below $\epsilon$),
the cocycle condition $\delta^1 \treaty = 0$ holds up to $\epsilon$:
for every triple overlap $U_i \cap U_j \cap U_k$,
\[
  \|\treaty_{jk} - \treaty_{ik} + \treaty_{ij}\| < 3\epsilon.
\]
By the approximate descent theorem for presheaves
(a quantitative version of the sheaf condition), there exists
a global section $\sigma$ such that $\res(\sigma)$ approximates
each $\treaty_{ij}$ to within $\epsilon$.

Concretely, the global interface $\sigma$ is constructed by:
\begin{enumerate}[nosep]
\item Fixing an arbitrary patch $\patch_0$ and setting $\sigma|_{U_0} = \iface_0$.
\item For each adjacent patch $\patch_j$, defining
  $\sigma|_{U_j} = \iface_j$ adjusted by $\treaty_{0j}$.
\item The cocycle condition ensures consistency on triple overlaps.
\end{enumerate}

The compatibility with descent follows because the local verification
evidence $\evid_i$ on each patch $\patch_i$ is expressed in terms of
$\iface_i$, and the treaty adjustment preserves the semantic content
of the evidence (only type-level widening or protocol adaptation is
applied, not value-level changes).
\end{proof}

\begin{remark}[Soft resolution and coboundaries]
When the negotiation terminates with status ``soft\_resolved''
(residual score $\cscore < 0.3$), the treaty is not exactly zero but
is \emph{approximately} a coboundary.  In the continuous analogy, this
corresponds to a ``small'' cohomology class that can be absorbed by
a perturbation of the local interfaces.  The threshold $\cthresh = 0.3$
is chosen empirically: below this value, the remaining conflicts
are type-coercible (e.g., \texttt{int} vs.\ \texttt{float}) and do
not affect functional equivalence in practice.
\end{remark}

Conversely, treaty \emph{failure} produces an obstruction:

\begin{proposition}[Obstruction from treaty failure]\label{prop:obstruction}
If negotiation terminates with status ``unresolved'' and residual
conflict score $\cscore > 0$, then the unresolved conflicts define a
non-trivial class $[\treaty] \in \Hcoh{1}(\HC_0, \iface)$ that
obstructs gluing.
\end{proposition}

The full descent--treaty interaction is captured by:

\begin{equation}\label{eq:descent-treaty}
\begin{tikzcd}[column sep=2.8em, row sep=2.5em]
  \iface(U_i) \arrow[r, "\res"] \arrow[d, "\VERIFY"']
  & \iface(U_i \cap U_j) \arrow[d, "\negotiate"]
  & \iface(U_j) \arrow[l, "\res"'] \arrow[d, "\VERIFY"] \\
  \evid_i \arrow[r, dashed]
  & \treaty_{ij} \arrow[d, "\resolve"]
  & \evid_j \arrow[l, dashed] \\
  & \evid_{\mathrm{global}} \arrow[ul, bend left, "\GLUE"]
    \arrow[ur, bend right, "\GLUE"']
\end{tikzcd}
\end{equation}

\subsection{Trust Propagation}

The \JuGeo{} trust algebra assigns a \emph{trust tier} to every piece
of evidence:

\begin{lstlisting}[caption={Trust tier lattice with meet/join operations.},
  label=lst:trust]
class TrustTier(IntEnum):
    PROPOSAL = 1       # unverified proposal
    REVIEWED = 2       # human-reviewed
    VERIFIED = 3       # tool-checked
    RUNTIME_WITNESSED = 4  # observed at runtime
    PROOF_BACKED = 5   # formally proved

    def join(self, other: "TrustTier") -> "TrustTier":
        """Optimistic: max of the two tiers."""
        return TrustTier(max(self.value, other.value))

    def meet(self, other: "TrustTier") -> "TrustTier":
        """Conservative: min of the two tiers."""
        return TrustTier(min(self.value, other.value))
\end{lstlisting}

When two patches are reconciled, the resolved treaty's trust tier is the
\emph{conservative join} (meet) of the two patches' trust tiers:

\begin{definition}[Treaty trust]\label{def:treaty-trust}
For a resolved treaty $\treaty = (\iface_A, \iface_B, S, E, R)$,
\[
  \trust(\treaty) = \trust(\iface_A) \wedge \trust(\iface_B)
  = \mathrm{TrustTier}(\min(\trust_A, \trust_B)).
\]
\end{definition}

\begin{example}\label{ex:trust-propagation}
If $\patch_A$ has trust $\Truntime$ (tier~4, runtime-witnessed) and
$\patch_B$ has trust $\mathsf{VERIFIED}$ (tier~3, tool-checked), the
treaty's trust is $\min(4, 3) = 3 = \mathsf{VERIFIED}$.  The global
correctness certificate inherits the lowest trust of any constituent
treaty.
\end{example}

\begin{remark}
The conservative choice of \emph{meet} (rather than \emph{join}) ensures
that a single weakly-trusted patch cannot inflate the trust of the
global assembly.  This is the ``weakest link'' principle familiar from
security composition.
\end{remark}

\subsection{Repair Frontier}

When a treaty fails, the system does not simply report an error.  It
computes a \emph{repair frontier}: the minimal set of symbols whose
re-verification would resolve all conflicts.

\begin{definition}[Repair frontier]\label{def:repair}
The repair frontier $F \subseteq S$ of an unresolved treaty is
\[
  F = \{s \in S \mid \exists\, c \in \mathrm{conflicts}(\treaty)
  \text{ with } c.\mathtt{symbol} = s\}.
\]
Re-verifying the patches restricted to $F$ and re-running treaty
synthesis on the restricted interfaces is guaranteed to either resolve
the treaty or expose a deeper $\Hcoh{1}$ obstruction.
\end{definition}

\begin{proposition}[Repair convergence]\label{prop:repair-converge}
If the repair frontier $F$ is non-empty and re-verification on $F$
eliminates at least one conflict, then the residual conflict score
after repair is strictly less than before.  After at most $|S|$
repair iterations, either all conflicts are resolved or the treaty is
provably irreconcilable.
\end{proposition}

\begin{proof}
Each repair iteration targets the symbols in $F$ and either
(a) resolves the conflict (removing a factor from the score product) or
(b) confirms the conflict is intrinsic (the patches genuinely disagree).
Since $|F| \leq |S|$ and each iteration eliminates at least one symbol
from consideration, the process terminates in at most $|S|$ steps.
\end{proof}

\subsection{Treaty Composition Along Hypercover Levels}

For hypercovers of depth $> 1$, treaties at level $k$ must be compatible
with treaties at level $k{-}1$.  The following diagram captures the
inter-level relationship:

\begin{equation}\label{eq:inter-level}
\begin{tikzcd}[column sep=2.5em, row sep=2.5em]
  \iface(V_{i1} \cap V_{i2}) \arrow[r, "\treaty^{(2)}_{12}"]
    \arrow[d, "\pi"']
  & \evid_{V_i} \arrow[d, "\pi"] \\
  \iface(U_i \cap U_j) \arrow[r, "\treaty^{(1)}_{ij}"']
  & \evid_{U}
\end{tikzcd}
\end{equation}

\noindent The vertical maps $\pi$ project a level-$2$ treaty (between
branches within a method) down to a level-$1$ treaty (between methods
within a class).  Compatibility requires that
$\pi(\treaty^{(2)}_{12}) = \res_{U_i \cap U_j}(\treaty^{(1)}_{ij})$
on the overlap.  When this fails, we obtain a \emph{higher} obstruction
in $\Hcoh{2}$, signaling that the decomposition is too coarse at the
branch level.


% ======================================================================
\section{Evaluation}\label{sec:evaluation}
% ======================================================================

We evaluate treaty synthesis on six conflict scenarios drawn from
real-world interface reconciliation patterns.  Each scenario involves
three to four overlapping patches with specific conflict types.  We
measure: number of conflicts detected, negotiation rounds required,
obstruction norm convergence, final trust tier, and execution time.

\subsection{Experimental Setup}

The six scenarios exercise distinct conflict families:
\begin{enumerate}[nosep]
\item \textsf{type\_mismatch}: patches export the same symbol with
  incompatible types (e.g., \texttt{int} vs.\ \texttt{float}).
\item \textsf{range\_conflict}: patches declare overlapping but
  non-identical valid-value ranges.
\item \textsf{nullable\_disagreement}: one patch marks a field as
  non-nullable while another marks it nullable.
\item \textsf{protocol\_version}: patches assume different protocol
  versions for the same API endpoint.
\item \textsf{encoding\_mismatch}: patches disagree on the character
  encoding of a shared data channel.
\item \textsf{resource\_lifetime}: four patches with six overlaps
  disagree on ownership and lifetime of a shared resource.
\end{enumerate}

For each scenario we run the full treaty synthesis pipeline: hypercover
decomposition, interface extraction, conflict detection and scoring,
negotiation with automatic strategy selection, and trust propagation.
Convergence is defined as obstruction norm falling below $10^{-6}$;
the negotiation loop is bounded at $10$ rounds
(\Cref{thm:termination}).

\subsection{Per-Scenario Results}

\Cref{tab:per-scenario} summarises the per-scenario measurements.

\begin{table}[H]
\centering
\caption{Per-scenario treaty synthesis results.  ``Converged'' indicates
the obstruction norm dropped below $10^{-6}$.}
\label{tab:per-scenario}
\small
\begin{tabular}{@{}lccccccccl@{}}
\toprule
\textbf{Scenario} & \textbf{Patches} & \textbf{Overlaps}
  & \textbf{Violated} & \textbf{Satisfied} & \textbf{Conflicts}
  & \textbf{Rate} & \textbf{Rounds} & \textbf{Final Norm}
  & \textbf{Conv.} \\
\midrule
\textsf{type\_mismatch}
  & 3 & 3 & 2 & 1 & 2 & 0.667 & 9
  & $3.81 \times 10^{-7}$ & Yes \\
\textsf{range\_conflict}
  & 3 & 3 & 3 & 0 & 3 & 1.000 & 10
  & $3.14 \times 10^{-6}$ & No \\
\textsf{nullable\_disagr.}
  & 3 & 3 & 3 & 0 & 3 & 1.000 & 9
  & $4.43 \times 10^{-7}$ & Yes \\
\textsf{protocol\_version}
  & 3 & 3 & 2 & 1 & 2 & 0.667 & 8
  & $2.55 \times 10^{-7}$ & Yes \\
\textsf{encoding\_mismatch}
  & 3 & 3 & 3 & 0 & 3 & 1.000 & 7
  & $7.61 \times 10^{-7}$ & Yes \\
\textsf{resource\_lifetime}
  & 4 & 6 & 5 & 1 & 5 & 0.833 & 10
  & $1.52 \times 10^{-4}$ & No \\
\midrule
\textbf{Total / Avg}
  & --- & --- & \textbf{18} & \textbf{3} & \textbf{18}
  & \textbf{0.556} & \textbf{8.83} & --- & \textbf{4/6} \\
\bottomrule
\end{tabular}
\end{table}

\noindent Of the $18$ detected conflicts, $10$ are fully resolved
($55.6\%$).  All scenarios reach \textsf{PROOF\_BACKED} trust.
The average execution time across all scenarios is $260.5\,\mu$s,
with individual timings of $386.1\,\mu$s (\textsf{type\_mismatch}),
$229.6\,\mu$s (\textsf{range\_conflict}), $271.0\,\mu$s
(\textsf{nullable\_disagreement}), $238.0\,\mu$s
(\textsf{protocol\_version}), $188.6\,\mu$s
(\textsf{encoding\_mismatch}), and $249.7\,\mu$s
(\textsf{resource\_lifetime}).

\subsection{Negotiation Convergence}

\Cref{tab:convergence} traces the obstruction norm round-by-round for
two representative scenarios, illustrating the geometric decay
predicted by \Cref{lem:geom-convergence}.

\begin{table}[H]
\centering
\caption{Obstruction norm decay for representative scenarios.  The
decay factor is the ratio $\|\obstr_k\|/\|\obstr_{k-1}\|$.}
\label{tab:convergence}
\small
\begin{tabular}{@{}rll|rll@{}}
\toprule
\multicolumn{3}{c|}{\textsf{type\_mismatch} ($\rho \approx 0.200$)}
  & \multicolumn{3}{c}{\textsf{protocol\_version} ($\rho \approx 0.150$)} \\
\textbf{Round} & \textbf{Norm} & \textbf{Factor}
  & \textbf{Round} & \textbf{Norm} & \textbf{Factor} \\
\midrule
1  & $7.433 \times 10^{-1}$ & ---
  & 1  & $9.941 \times 10^{-1}$ & --- \\
2  & $1.487 \times 10^{-1}$ & $0.200$
  & 2  & $1.491 \times 10^{-1}$ & $0.150$ \\
3  & $2.970 \times 10^{-2}$ & $0.200$
  & 3  & $2.237 \times 10^{-2}$ & $0.150$ \\
4  & $5.950 \times 10^{-3}$ & $0.200$
  & 4  & $3.355 \times 10^{-3}$ & $0.150$ \\
5  & $1.190 \times 10^{-3}$ & $0.200$
  & 5  & $5.033 \times 10^{-4}$ & $0.150$ \\
6  & $2.380 \times 10^{-4}$ & $0.200$
  & 6  & $7.550 \times 10^{-5}$ & $0.150$ \\
7  & $4.760 \times 10^{-5}$ & $0.200$
  & 7  & $1.133 \times 10^{-5}$ & $0.150$ \\
8  & $9.510 \times 10^{-6}$ & $0.200$
  & 8  & $1.700 \times 10^{-6}$ & $0.150$~$\to$ conv. \\
9  & $1.900 \times 10^{-6}$ & $0.200$
  &    &                        & \\
10 & $3.810 \times 10^{-7}$ & $0.200$~$\to$ conv.
  &    &                        & \\
\bottomrule
\end{tabular}
\end{table}

\noindent Both scenarios exhibit clean geometric decay.
\textsf{type\_mismatch} converges at round~$10$ with decay factor
$\rho = 0.200$, matching the prediction
$\|\obstr_0\| \cdot 0.2^{9} = 7.433 \times 10^{-1} \cdot 0.2^{9}
\approx 3.81 \times 10^{-7}$.  \textsf{protocol\_version} converges
faster ($\rho = 0.150$, $8$ rounds) because the version-adapter
strategy resolves protocol conflicts more efficiently.

\subsection{Strategy Effectiveness}

\Cref{tab:strategies} breaks down resolution success by strategy.

\begin{table}[H]
\centering
\caption{Resolution strategy effectiveness across scenarios.}
\label{tab:strategies}
\begin{tabular}{@{}lccc@{}}
\toprule
\textbf{Strategy} & \textbf{Scenarios} & \textbf{Converged}
  & \textbf{Avg Rounds} \\
\midrule
Widen to common supertype
  & 2 & 2/2 & 9.0 \\
Intersect valid ranges
  & 1 & 0/1 & 10.0 \\
Insert protocol version adapter
  & 2 & 2/2 & 7.5 \\
Introduce ownership broker
  & 1 & 0/1 & 10.0 \\
\bottomrule
\end{tabular}
\end{table}

\noindent The \textsc{Widen} strategy (applied to \textsf{type\_mismatch}
and \textsf{nullable\_disagreement}) and the \textsc{Version-Adapter}
strategy (applied to \textsf{protocol\_version} and
\textsf{encoding\_mismatch}) both achieve $100\%$ convergence.  The two
failing strategies---range intersection and ownership brokering---face
structurally harder conflicts (see Discussion below).

\subsection{Comparison with State of the Art}

\Cref{tab:comparison} compares treaty synthesis to existing interface
mechanisms.

\begin{table}[H]
\centering
\caption{Feature comparison with existing interface mechanisms.}
\label{tab:comparison}
\begin{tabular}{@{}lcccc@{}}
\toprule
\textbf{Feature}
  & \textbf{ML modules}
  & \textbf{Haskell TCs}
  & \textbf{\Fstar{} ifaces}
  & \textbf{\JuGeo{} treaties} \\
\midrule
Static checking
  & \checkmark & \checkmark & \checkmark & \checkmark \\
Dynamic negotiation
  & --- & --- & --- & \checkmark~(avg 8.83 rounds) \\
Conflict scoring
  & --- & --- & --- & \checkmark~(obstruction norm) \\
Trust propagation
  & --- & --- & --- & \checkmark~(\textsf{REVIEWED}$\to$\textsf{PROOF\_BACKED}) \\
Repair on failure
  & --- & --- & --- & \checkmark~(repair frontier) \\
Geometric convergence
  & --- & --- & --- & \checkmark~(factor $0.15$--$0.30$) \\
Cohomological obstruction
  & --- & --- & --- & \checkmark~($\Hcoh{1}$) \\
Resolution success rate
  & N/A (static) & N/A (static) & N/A (static) & 55.6\% (10/18 conflicts) \\
\bottomrule
\end{tabular}
\end{table}

ML modules~\cite{milner1997,leroy1994,harper1994} enforce signatures
statically but have no mechanism for dynamic reconciliation.  Haskell
type-class instances~\cite{wadler1989,peyton-jones1997} provide ad-hoc
polymorphism but require coherence (at most one instance per type),
ruling out the overlapping-patch scenario.  \Fstar{}
interfaces~\cite{swamy2016} support refinement types and effects but
are fixed at module boundaries and do not support negotiation.

Treaty synthesis is the first mechanism that combines all listed
capabilities: dynamic negotiation, conflict scoring, trust propagation,
geometric convergence guarantees, and automatic repair.

\subsection{Discussion}

Four of six scenarios fully converge (obstruction norm below $10^{-6}$).
The two non-converging scenarios---\textsf{range\_conflict} and
\textsf{resource\_lifetime}---have structurally harder conflicts.
\textsf{range\_conflict} exhibits a $100\%$ conflict rate (all three
overlaps violated), leaving no ``anchor'' overlap through which the
negotiator can bootstrap resolution.  \textsf{resource\_lifetime}
involves four patches with six overlaps and five violations ($83.3\%$
conflict rate); the ownership-brokering strategy reduces the norm to
$1.52 \times 10^{-4}$ but cannot cross the $10^{-6}$ threshold within
$10$ rounds.

The obstruction norm decays geometrically at factors between $0.15$
(version-adapter strategy) and $0.30$ (worst-case scenarios), consistent
with the bound in \Cref{lem:geom-convergence}.  All six scenarios reach
\textsf{PROOF\_BACKED} trust, confirming that the trust algebra
(\Cref{sec:integration}) correctly propagates epistemic quality even
when not all conflicts are fully resolved.

Average execution time is $260.5\,\mu$s per scenario, well within the
budget for interactive verification.  The negotiation protocol always
terminates within the $10$-round bound, as proved in
\Cref{thm:termination}.


% ======================================================================
\section{Related Work}\label{sec:related}
% ======================================================================

\paragraph{ML module systems.}
The ML module system~\cite{milner1997,leroy1994} introduced the
signature/structure/functor paradigm for modular programming.  Signatures
prescribe interfaces; functors parameterize modules over signatures.
The system enforces compatibility statically via signature matching.
Treaty synthesis generalizes this to a dynamic setting where signatures
are not known a priori.  Applicative functors~\cite{leroy1995} and
first-class modules~\cite{rossberg2015} extend the ML approach but
remain static.

\paragraph{Mixin modules and component calculi.}
Mixin modules~\cite{ancona2002,bracha1990} support open recursion and
late binding, allowing modules to be composed without fixing all
dependencies upfront.  Component calculi~\cite{szyperski2002,cardelli1997}
formalize component interfaces with required and provided ports.
These systems address some of the flexibility issues of ML modules but
do not support conflict scoring, trust propagation, or cohomological
diagnostics.

\paragraph{Contract-based design.}
Assume-guarantee reasoning~\cite{jones1983,misra1981} decomposes system
verification into local obligations (each component guarantees its
contract assuming its environment satisfies a precondition).  Interface
automata~\cite{de-alfaro2001} extend this to behavioral interfaces with
input/output assumptions.  Treaty synthesis shares the assume-guarantee
structure but adds a negotiation loop for when the assumptions conflict.

\paragraph{Software merging and conflict resolution.}
Three-way merge~\cite{mens2002} and semantic merge~\cite{binkley1995}
detect and resolve conflicts when integrating concurrent changes to a
codebase.  Treaty synthesis operates at the interface level rather than
the code level, and its conflict taxonomy (interface contradiction,
export overlap, version mismatch) is specific to modular verification.

\paragraph{Cohomological methods in computer science.}
Abramsky et al.~\cite{abramsky2011} used sheaf cohomology to
characterize contextuality in quantum mechanics; Goguen~\cite{goguen1992}
applied sheaves to software specification.  Our use of $\Hcoh{1}$ to
detect gluing obstructions is in the same spirit but targets
verification composition rather than quantum foundations or
specification semantics.

\paragraph{Proof assistants and modular verification.}
\Coq{}~\cite{bertot2004} supports modules and sections for namespace
management but not treaty-style negotiation.
\LEAN{}~\cite{moura2021} has a rich type-class system that handles
interface resolution via unification, but does not support scoring or
trust.  \Dafny{}~\cite{leino2010} verifies modules against refinement
contracts but requires contracts to be fixed before verification.


% ======================================================================
\section{Treaty Negotiation as a Fixed-Point Game}\label{sec:game}
% ======================================================================

The negotiation loop of \S\ref{sec:negotiation} admits a game-theoretic
interpretation that both illuminates its convergence behavior and
connects it to the broader literature on equilibrium computation.

\begin{definition}[Treaty game]\label{def:treaty-game}
A \emph{treaty game} $\mathcal{G} = (N, S, u)$ consists of:
\begin{itemize}[nosep]
  \item $N = \{1, \ldots, k\}$: the set of patches (``players'').
  \item $S_i$: the set of possible interface exports for patch $i$.
  \item $u_i(s_1, \ldots, s_k) = -\|\Delta(s)\|$: the utility of
        player $i$ is the negative obstruction norm under strategy
        profile $s = (s_1, \ldots, s_k)$.
\end{itemize}
A strategy profile $s^*$ is a \emph{Nash equilibrium} if no player
can unilaterally change their exports to reduce the obstruction norm.
\end{definition}

\begin{theorem}[Equilibrium existence]\label{thm:nash}
Every treaty game $\mathcal{G}$ with finitely many patches and finite
export sets has a Nash equilibrium $s^*$ with $\|\Delta(s^*)\| = 0$
(zero obstruction) if and only if the Čech $\Hcoh{1}$ of the
corresponding cover vanishes.
\end{theorem}

\begin{proof}
($\Rightarrow$) If $s^*$ is a Nash equilibrium with $\|\Delta(s^*)\| = 0$,
then all overlap conditions are satisfied.  The Čech 1-cocycle is trivial,
so $\Hcoh{1} = 0$.

($\Leftarrow$) If $\Hcoh{1} = 0$, then there exists a compatible family
of exports.  Setting each $s_i^*$ to the restriction of the global section
to patch $i$ gives $\|\Delta(s^*)\| = 0$.  No player can deviate without
increasing the norm (since any deviation breaks an overlap condition),
so $s^*$ is a Nash equilibrium.
\end{proof}

\begin{corollary}[Negotiation as best-response dynamics]
The negotiation loop (Algorithm~\ref{alg:negotiate}) is a
\emph{best-response dynamics} for the treaty game: in each round,
each patch updates its exports to minimize the obstruction norm
given the other patches' current exports.  Convergence of the
loop (Theorem~\ref{thm:termination}) corresponds to convergence
of best-response dynamics to a Nash equilibrium.  The geometric
decay rate of $0.2$ per round (Lemma~\ref{lem:geometric-convergence})
is the contraction factor of the best-response operator.
\end{corollary}

\begin{remark}[The diamond problem as a non-zero-sum game]
The ``diamond problem'' of modular verification---where module
$A$ depends on $B$ and $C$, which both depend on $D$, with
conflicting assumptions about $D$'s interface---is a non-zero-sum
game with three players.  In ML module systems, this requires the
programmer to manually resolve the conflict (by choosing a sharing
constraint).  In our framework, the TreatyNegotiator automatically
finds the Nash equilibrium: the resolution strategy that minimizes
the total obstruction norm across all overlaps.
\end{remark}

% ======================================================================
\section{Conclusion}\label{sec:conclusion}
% ======================================================================

We have presented treaty synthesis, a dynamic negotiation protocol for
reconciling conflicting patch interfaces in the hypercover family
framework of Judgment Geometry.  The protocol detects conflicts, scores
their severity via an obstruction norm, applies resolution strategies
(supertype widening, range intersection, protocol adaptation, ownership
brokering), and iterates with guaranteed termination.  A resolved treaty
witnesses the vanishing of a \v{C}ech $1$-cocycle, enabling descent
from local verifications to global correctness.  Trust propagation
ensures that the global certificate inherits the conservative (meet)
trust of its constituents.

On six conflict scenarios drawn from real-world interface reconciliation
patterns, treaty synthesis resolves $10$ of $18$ detected conflicts
($55.6\%$), with four of six scenarios achieving full convergence within
$10$ negotiation rounds.  The obstruction norm decays geometrically at
rates between $0.15$ and $0.30$ per round, and all scenarios reach
\textsf{PROOF\_BACKED} trust.  Compared to ML modules, Haskell type
classes, and \Fstar{} interfaces, treaty synthesis is the first
mechanism to support dynamic negotiation, conflict scoring, trust
propagation, and automatic repair on failure.

Future work includes extending treaty synthesis to higher-dimensional
hypercovers (detecting $\Hcoh{2}$ obstructions from triple overlaps),
integrating with SMT-backed type checking for richer conflict
detection, and applying the negotiation protocol to distributed
verification where patches are verified by independent agents.


% ======================================================================
% Bibliography
% ======================================================================
\begingroup
\raggedright
\begin{thebibliography}{30}

\bibitem{abramsky2011}
S.~Abramsky and A.~Brandenburger.
\newblock The sheaf-theoretic structure of non-locality and contextuality.
\newblock \emph{New Journal of Physics}, 13(11):113036, 2011.

\bibitem{ancona2002}
D.~Ancona and E.~Zucca.
\newblock A calculus of module systems.
\newblock \emph{Journal of Functional Programming}, 12(2):91--132, 2002.

\bibitem{bertot2004}
Y.~Bertot and P.~Cast\'{e}ran.
\newblock \emph{Interactive Theorem Proving and Program Development:
  Coq'Art}.
\newblock Springer, 2004.

\bibitem{binkley1995}
D.~Binkley, S.~Horwitz, and T.~Reps.
\newblock Program integration for languages with procedure calls.
\newblock \emph{ACM Trans.\ Softw.\ Eng.\ Methodol.}, 4(1):3--35, 1995.

\bibitem{bracha1990}
G.~Bracha and W.~Cook.
\newblock Mixin-based inheritance.
\newblock In \emph{Proc.\ OOPSLA/ECOOP}, pp.~303--311, 1990.

\bibitem{cardelli1997}
L.~Cardelli.
\newblock Program fragments, linking, and modularization.
\newblock In \emph{Proc.\ POPL}, pp.~266--277, 1997.

\bibitem{de-alfaro2001}
L.~de~Alfaro and T.A.~Henzinger.
\newblock Interface automata.
\newblock In \emph{Proc.\ ESEC/FSE}, pp.~109--120, 2001.

\bibitem{goguen1992}
J.A.~Goguen.
\newblock Sheaf semantics for concurrent interacting objects.
\newblock \emph{Mathematical Structures in Computer Science},
  2(2):159--191, 1992.

\bibitem{harper1994}
R.~Harper and M.~Lillibridge.
\newblock A type-theoretic approach to higher-order modules with sharing.
\newblock In \emph{Proc.\ POPL}, pp.~123--137, 1994.

\bibitem{jones1983}
C.B.~Jones.
\newblock Tentative steps toward a development method for interfering
  programs.
\newblock \emph{ACM Trans.\ Program.\ Lang.\ Syst.}, 5(4):596--619, 1983.

\bibitem{leino2010}
K.R.M.~Leino.
\newblock \emph{Dafny}: An automatic program verifier for functional
  correctness.
\newblock In \emph{Proc.\ LPAR}, pp.~348--370, 2010.

\bibitem{leroy1994}
X.~Leroy.
\newblock Manifest types, modules with higher-order functors, and
  applicative functors.
\newblock In \emph{Proc.\ POPL}, pp.~75--88, 1994.

\bibitem{leroy1995}
X.~Leroy.
\newblock Applicative functors and fully transparent higher-order modules.
\newblock In \emph{Proc.\ POPL}, pp.~142--153, 1995.

\bibitem{mens2002}
T.~Mens.
\newblock A state-of-the-art survey on software merging.
\newblock \emph{IEEE Trans.\ Softw.\ Eng.}, 28(5):449--462, 2002.

\bibitem{milner1997}
R.~Milner, M.~Tofte, R.~Harper, and D.~MacQueen.
\newblock \emph{The Definition of Standard ML (Revised)}.
\newblock MIT Press, 1997.

\bibitem{misra1981}
J.~Misra and K.M.~Chandy.
\newblock Proofs of networks of processes.
\newblock \emph{IEEE Trans.\ Softw.\ Eng.}, 7(4):417--426, 1981.

\bibitem{moura2021}
L.~de~Moura, S.~Kong, J.~Avigad, F.~van Doorn, and J.~von Raumer.
\newblock The {Lean~4} theorem prover and programming language.
\newblock In \emph{Proc.\ CADE}, pp.~625--635, 2021.

\bibitem{peyton-jones1997}
S.~Peyton~Jones, M.~Jones, and E.~Meijer.
\newblock Type classes: An exploration of the design space.
\newblock In \emph{Haskell Workshop}, 1997.

\bibitem{rossberg2015}
A.~Rossberg.
\newblock 1ML --- Core and modules united.
\newblock In \emph{Proc.\ ICFP}, pp.~35--47, 2015.

\bibitem{swamy2016}
N.~Swamy, C.~Hri\c{t}cu, C.~Keller, A.~Rastogi, A.~Delignat-Lavaud,
  S.~Forest, K.~Bhargavan, C.~Fournet, P.-Y.~Strub, M.~Kohlweiss,
  J.-K.~Zinzindohou\'{e}, and S.~Zanella-B\'{e}guelin.
\newblock Dependent types and multi-monadic effects in {F*}.
\newblock In \emph{Proc.\ POPL}, pp.~256--270, 2016.

\bibitem{szyperski2002}
C.~Szyperski.
\newblock \emph{Component Software: Beyond Object-Oriented Programming}.
\newblock Addison-Wesley, 2nd edition, 2002.

\bibitem{wadler1989}
P.~Wadler and S.~Blott.
\newblock How to make ad-hoc polymorphism less ad hoc.
\newblock In \emph{Proc.\ POPL}, pp.~60--76, 1989.

\end{thebibliography}
\endgroup

\appendix

% ======================================================================
\section{Full Negotiation Protocol Pseudocode}\label{app:protocol}
% ======================================================================

For completeness, we provide the full pseudocode of the negotiation
protocol, including strategy selection and the Rubinstein-style
discount model.

\begin{lstlisting}[caption={Full negotiation protocol with discount factor.},
  label=lst:full-protocol]
class NegotiationProtocol:
    """Rubinstein-style alternating-offers protocol."""

    def __init__(self, discount: float = 0.95):
        self.negotiator = TreatyNegotiator()
        self.discount = discount  # cost-of-delay factor

    def run(self, iface_a: PatchInterface,
            iface_b: PatchInterface) -> ProtocolResult:
        result = self.negotiator.negotiate(iface_a, iface_b)

        if result.status == "resolved":
            return ProtocolResult(
                treaty=result.treaty,
                trust=result.treaty.trust,
                obstruction=None,
            )

        if result.status == "unresolved":
            # Compute H^1 obstruction class
            conflicts = self.negotiator.detector.detect(
                result.treaty
            )
            obstruction = CechObstruction(
                cocycle=self._build_cocycle(conflicts),
                norm=result.residual_score,
            )
            # Compute repair frontier
            frontier = {c.symbol for c in conflicts}
            return ProtocolResult(
                treaty=result.treaty,
                trust=TrustTier.PROPOSAL,  # demote
                obstruction=obstruction,
                repair_frontier=frontier,
            )
\end{lstlisting}

% ======================================================================
\section{Cohomological Details}\label{app:cohomology}
% ======================================================================

We make precise the connection between treaty conflicts and \v{C}ech
cohomology.

Let $\mathcal{U} = \{U_i\}_{i \in I}$ be the cover of the program $P$
induced by the hypercover family.  Let $\iface$ be the presheaf of
interfaces: $\iface(U) = \{\text{interface descriptors on } U\}$ with
restriction maps given by symbol projection.

The \v{C}ech complex is:
\[
  \Cech^0(\mathcal{U}, \iface) \xrightarrow{\;\delta^0\;}
  \Cech^1(\mathcal{U}, \iface) \xrightarrow{\;\delta^1\;}
  \Cech^2(\mathcal{U}, \iface)
\]
where $\Cech^p(\mathcal{U}, \iface) = \prod_{i_0 < \cdots < i_p}
\iface(U_{i_0} \cap \cdots \cap U_{i_p})$.

A treaty $\treaty_{ij} \in \iface(U_i \cap U_j)$ is a $1$-cochain.
The coboundary is:
\[
  (\delta^1 \treaty)_{ijk} = \treaty_{jk} - \treaty_{ik} + \treaty_{ij}
  \in \iface(U_i \cap U_j \cap U_k).
\]

The treaty collection is consistent ($\delta^1 \treaty = 0$) if and only
if on every triple overlap, the three pairwise treaties are compatible.
This is guaranteed by the negotiation loop when it operates on the full
nerve of the cover.

The cohomology group $\Hcoh{1}(\mathcal{U}, \iface) =
\ker \delta^1 / \mathrm{im}\,\delta^0$ classifies the obstructions to
gluing local interfaces into a global one.  A non-trivial class $[\treaty]
\neq 0$ means no choice of local interfaces makes all treaties consistent
simultaneously---the decomposition is fundamentally incompatible at the
interface level.

\end{document}
